\section*{Source-Grounded Synthesis Across the Propers}

\subsection*{From provision to the giver}

The correlated reading and Gospel begin with hunger answered by provision,
but their shared movement does not stop at provision. Exodus makes manna a
gift whose unfamiliarity provokes a question; Psalm 78 turns the event into
received memory; John receives the manna tradition while correcting both its
giver and its horizon. Moses did not originate the heavenly bread. The Father
gives the true bread, and the one speaking is himself the bread of life.

\subsection*{Gift that elicits a new manner of life}

Ephesians remains semi-continuous, yet it prevents the Sunday from becoming a
theory of religious consumption. The crowd asks what work it should perform;
Jesus answers with belief in the one sent. The apostolic reading describes
learning Christ as putting off the old person, being renewed in mind, and
putting on the new. Grace precedes this response, as manna precedes Israel's
recognition, but grace also forms conduct.

\subsection*{Bread, word, and heavenly gift}

The acclamation denies that bread alone constitutes life. John does not
therefore make bread irrelevant: Jesus has fed the crowd, and his
self-identification uses hunger to reveal a greater gift. The first Communion
option receives Wisdom's heavenly-bread image; the second repeats John's
bread-of-life saying. The final prayer describes reception as a heavenly gift
ordered toward eternal redemption. These liturgical statements permit a
eucharistic horizon for the whole, while the commentary avoids claiming that
every use of bread in the readings is already an exhaustive sacramental
definition.

\subsection*{The alternative endings remain alternatives}

At Communion, John~6:35 creates the shortest direct return to the Gospel:
Christ names himself as bread and promises an end to hunger and thirst for
those who come and believe. Wisdom~16:20 instead returns to heavenly bread as
delight adapted to the receiver. Both can lead into the final petition for
protection and eternal redemption, but they do so differently. A preparation
guide may describe both printed options;
it may not narrate them as one enacted sequence.

\begin{fourstable}{Movement}{Primary anchors}{Checked reception}{Controlling limit}
Gift and question & Exod.~16; Ps.~78; John~6 & Augustine; Chrysostom & Bodily hunger is real; the giver exceeds the gift.\\
Belief and renewal & John~6; Eph.~4 & Chrysostom on John and Ephesians & The apostolic course remains semi-continuous.\\
Bread and heavenly gift & Matt.~4; John~6; Communion A/B; final prayer & Augustine; Ambrose & Typology does not erase textual distinctions.\\
Wisdom or direct return & Communion A/B; final prayer & Wisdom and Johannine reception & The alternatives are not combined.\\
\end{fourstable}

\clearpage
